% 1-3 题 图 1-2：一系列光滑斜面（顶端与底端水平距离同为 l，倾角 φ 连续取值）
\begin{tikzpicture}[>=Stealth, thick]

    % 定义汇聚点坐标
    \coordinate (A) at (4,0);
    
    % 1. 绘制水平地面线
    \draw (-1,0) -- (5,0);
    
    % 绘制地面的斜线阴影 (表示刚性固定面)
    \foreach \x in {-0.8,-0.4,...,4.8} {
        \draw (\x,0) -- (\x-0.2,-0.2);
    }

    % 2. 绘制竖直虚线
    \draw[dashed] (0,0) -- (0,4.5);

    % 3. 绘制三根连接到地面的斜线
    \draw (0,1.2) -- (A);
    \draw (0,2.6) -- (A);
    \draw (0,4.0) -- (A);

    % 4. 绘制底部的长度标注 l
    % 向下延伸辅助虚线
    \draw[dashed, thin] (0,0) -- (0,-0.8);
    \draw[dashed, thin] (4,0) -- (4,-0.8);
    % 绘制双向箭头及标签
    \draw[<->] (0,-0.6) -- (4,-0.6) node[midway, fill=white, inner sep=2pt] {$l$};

    % 5. 绘制角度 \phi
    % 计算角度: arctan(1.2 / 4) ≈ 16.7°
    % 圆弧中心在 (4,0)，从左侧水平面 (180°) 画到第一根斜线处 (180 - 16.7 = 163.3°)
    \draw (3.2,0) arc[start angle=180, end angle=163.3, radius=0.8];
    \node at (2.9, 0.25) {$\phi$};

\end{tikzpicture}
